\section{Desarrollo}

\subsection{Uso de Alembic a lo largo del desarrollo}\label{sec:alembic_uso}

El desarrollo del backend no siguió un esquema rígido en el que el modelo de base de datos se diseña una sola vez y permanece intacto. Por el contrario, a medida que se avanzó en la implementación de los servicios y, más tarde, en la construcción del frontend, surgieron necesidades que requirieron ajustes en el esquema. Alembic fue la herramienta que permitió gestionar estos cambios de forma controlada y trazable, generando un script de migración por cada modificación relevante.

A lo largo del proyecto se generaron un total de 8 migraciones versionadas, nombradas 
bajo la convención \texttt{Parche X.Y}, que documentan la evolución del esquema desde 
su creación hasta su forma final. La migración inicial estableció todas las tablas del 
sistema. Las siguientes migraciones atendieron ajustes puntuales que surgieron durante 
el desarrollo: la ampliación del campo \texttt{password} para compatibilidad con 
\texttt{pwdlib}, la eliminación del campo \texttt{comment} de \texttt{device\_status\_logs} 
al redefinir la responsabilidad del módulo, y el registro explícito de \texttt{UserRole} 
como modelo ORM para que SQLAlchemy gestionara correctamente la tabla intermedia. 
Finalmente, los parches 1.5 a 1.8 resolvieron una serie de problemas con la propagación 
de eliminaciones en cascada (\texttt{ON DELETE CASCADE}), detectados al implementar el 
frontend: al eliminar un \texttt{product}, sus \texttt{devices} no se eliminaban; y al 
eliminar un \texttt{device}, sus \texttt{loan\_request\_items} quedaban huérfanos. Las 
relaciones corregidas fueron \texttt{devices} $\rightarrow$ \texttt{products}, 
\texttt{loan\_request\_items} $\rightarrow$ \texttt{devices}, y 
\texttt{loan\_request\_items} $\rightarrow$ \texttt{loan\_requests}.

Este historial de migraciones ilustra de forma concreta cómo Alembic actúa como red de seguridad durante el desarrollo iterativo: cada cambio queda documentado, es reversible mediante \texttt{downgrade}, y puede aplicarse en cualquier entorno con un único comando, garantizando coherencia entre el modelo Python y el esquema real de la base de datos.

% -------------------------------------------------------
\subsection{Módulo de Usuarios}\label{sec:backend_usuarios}

El módulo de \texttt{users} concentra la gestión de identidad y datos personales de quienes interactúan con el sistema. Su implementación consta de un repositorio, un conjunto de DTOs y un controlador que expone la API REST.

\subsubsection*{DTOs}

Se definen tres DTOs para el modelo \texttt{User}:

\begin{itemize}
	\item \texttt{UserCreateDTO}: utilizado al crear un usuario; excluye \texttt{id} y \texttt{created\_at}, que se generan automáticamente.
	\item \texttt{UserReadDTO}: utilizado en todas las respuestas; excluye \texttt{password} para evitar exponer credenciales.
	\item \texttt{UserUpdateDTO}: utilizado en actualizaciones parciales; todos los campos son opcionales (\texttt{partial=True}).
	\item \texttt{UserLoginDTO}: utilizado exclusivamente por el servicio \texttt{auth} para el flujo de login; expone únicamente \texttt{username} y \texttt{password}.
\end{itemize}

\subsubsection*{Repositorio}

\texttt{UserRepository} extiende \texttt{SQLAlchemySyncRepository} e incorpora tres métodos personalizados:

\begin{itemize}
	\item \texttt{get\_my\_user(username)}: recupera un usuario por su \texttt{username} con todas sus relaciones cargadas mediante \texttt{selectinload} (\texttt{roles}, \texttt{loan\_requests}, \texttt{status\_logs}). Es utilizado por el sistema de autenticación para resolver el usuario a partir del token JWT.
	
	\item \texttt{add\_with\_existing\_roles(roles\_repo, user)}: crea un usuario hasheando su contraseña con \texttt{pwdlib} y asociándolo a roles ya existentes en la base de datos, consultados a través de \texttt{RoleRepository}. Este método personalizado es necesario porque la creación de un usuario involucra dos tablas (\texttt{users} y \texttt{user\_roles}) y requiere cifrar la contraseña antes de persistir.
	
	\item \texttt{update\_with\_existing\_roles(roles\_repo, user\_id, user\_data)}: actualiza los campos de un usuario y, si el payload incluye roles, resuelve los objetos \texttt{Role} correspondientes antes de aplicar el cambio. Permite actualizaciones parciales respetando el contrato de \texttt{UserUpdateDTO}.
	
	\item \texttt{check\_password(username, password)}: verifica la contraseña de un usuario comparando el texto plano contra el hash almacenado mediante \texttt{pwdlib}. Utilizado por el controlador de \texttt{auth} durante el login.
\end{itemize}

Los métodos genéricos \texttt{list}, \texttt{get\_one} y \texttt{delete} provienen directamente de \texttt{SQLAlchemySyncRepository} y no requieren implementación adicional.

\subsubsection*{Controlador y guards}

\texttt{UserController} expone la ruta base \texttt{/users} e implementa los siguientes endpoints:

\begin{itemize}
	\item \texttt{GET /users/me}: retorna el usuario autenticado extrayéndolo del objeto \texttt{request}. No requiere parámetros adicionales.
	\item \texttt{GET /users/}: lista todos los usuarios mediante \texttt{list()}.
	\item \texttt{GET /users/\{user\_id\}}: obtiene un usuario por ID mediante \texttt{get\_one()}; retorna \texttt{404} si no existe.
	\item \texttt{POST /users/}: crea un usuario con roles; protegido por \texttt{admin\_user\_guard}. Delega en \texttt{add\_with\_existing\_roles}.
	\item \texttt{PATCH /users/\{user\_id\}}: actualiza un usuario; protegido por \texttt{admin\_user\_guard}. Delega en \texttt{update\_with\_existing\_roles}.
	\item \texttt{DELETE /users/\{user\_id\}}: elimina un usuario; protegido por \texttt{admin\_user\_guard}.
\end{itemize}

Se definen dos guards en este módulo:

\begin{itemize}
	\item \texttt{require\_active}: aplicado a nivel de controlador; rechaza con \texttt{403} cualquier petición de un usuario cuya cuenta esté inactiva (\texttt{is\_active = False}). Garantiza que ningún usuario deshabilitado pueda operar el sistema.
	\item \texttt{admin\_user\_guard}: aplicado a los endpoints de escritura (\texttt{POST}, \texttt{PATCH}, \texttt{DELETE}); rechaza con \texttt{403} si el usuario autenticado no tiene \texttt{is\_admin = True}. Ambas guards reciben \texttt{ASGIConnection} para acceder al usuario ya resuelto por el middleware JWT.
\end{itemize}

% -------------------------------------------------------
\subsection{Módulo de Autenticación}\label{sec:backend_auth}

La autenticación se implementa a través del servicio \texttt{auth}, cuyo controlador \texttt{AuthController} expone un único endpoint:

\begin{itemize}
	\item \texttt{POST /auth/login}: recibe las credenciales del usuario codificadas como \texttt{application/x-www-form-urlencoded} mediante \texttt{UserLoginDTO}. Verifica la existencia del usuario con \texttt{get\_one\_or\_none} y la validez de la contraseña con \texttt{check\_password}; si alguna falla, responde con \texttt{401}. En caso de éxito, delega en \texttt{oauth2\_auth.login} para emitir el token JWT.
\end{itemize}

La configuración de seguridad centralizada reside en \texttt{security.py}, donde se instancia \texttt{OAuth2PasswordBearerAuth} con una función \texttt{retrieve\_user\_handler} que, dado un token, abre sesión en la base de datos y recupera el usuario mediante \texttt{get\_my\_user}. Este objeto \texttt{oauth2\_auth} se registra como middleware de la aplicación Litestar, de modo que todas las rutas protegidas reciben automáticamente el usuario autenticado en \texttt{request.user}, disponible también en las guards a través de \texttt{ASGIConnection}.

% -------------------------------------------------------
\subsection{Módulo de Catálogo}\label{sec:backend_catalogo}

El catálogo agrupa cinco entidades de referencia: \texttt{brands}, \texttt{models}, \texttt{categories}, \texttt{statuses} y \texttt{ubications}, además de \texttt{roles}, que fue integrado aquí por compartir la misma naturaleza de dato de referencia con métodos acotados. Todos estos recursos siguen una estructura homogénea y predecible.

Cada entidad del catálogo posee su propio par de DTOs (\texttt{ReadDTO}, \texttt{CreateDTO}, \texttt{UpdateDTO}), su repositorio basado en \texttt{SQLAlchemySyncRepository}, y un controlador con los cuatro endpoints genéricos: \texttt{GET /}, \texttt{GET /\{id\}}, \texttt{POST /}, \texttt{PATCH /\{id\}} y \texttt{DELETE /\{id\}}. Los DTOs excluyen siempre las relaciones inversas para evitar serialización circular (por ejemplo, \texttt{BrandReadDTO} excluye \texttt{products}).

Todos los métodos utilizados son genéricos y provienen directamente de \texttt{SQLAlchemySyncRepository}: \texttt{list}, \texttt{get\_one}, \texttt{add}, \texttt{get\_and\_update} y \texttt{delete}. No se implementó ningún método personalizado en este módulo, ya que la simplicidad de sus entidades no lo requiere.

Una excepción destacable es el controlador de \texttt{roles}, que define la constante \texttt{PROTECTED\_ROLES = \{`usuario'\}}  y verifica en el endpoint \texttt{DELETE} que el rol a eliminar no pertenezca a dicho conjunto. De ser así, responde con \texttt{409 Conflict} para proteger roles del sistema que son necesarios para el funcionamiento correcto de la plataforma.

% -------------------------------------------------------
\subsection{Módulo de Productos}\label{sec:backend_productos}

El módulo de \texttt{products} gestiona los tipos de ítems disponibles en el inventario. Su estructura sigue el patrón estándar del proyecto con tres DTOs, un repositorio y un controlador.

\subsubsection*{DTOs}

\begin{itemize}
	\item \texttt{ProductReadDTO}: excluye \texttt{devices} para evitar serialización de todas las unidades físicas al listar productos.
	\item \texttt{ProductCreateDTO}: excluye \texttt{id}, \texttt{created\_at} y \texttt{devices}.
	\item \texttt{ProductUpdateDTO}: igual que \texttt{CreateDTO} pero con \texttt{partial=True}.
\end{itemize}

\subsubsection*{Repositorio y método personalizado de actualización}

\texttt{ProductRepository} utiliza exclusivamente métodos genéricos para \texttt{list}, \texttt{get\_one} y \texttt{delete}. Sin embargo, el endpoint de actualización implementa lógica personalizada directamente en el controlador, dado que \texttt{as\_builtins()} con \texttt{partial=True} omite los campos con valor \texttt{None}, lo que impediría limpiar relaciones opcionales como \texttt{brand\_id}, \texttt{model\_id} o \texttt{category\_id}. Para resolverlo, el controlador aplica estos tres campos directamente sobre la instancia con \texttt{setattr}, forzando a SQLAlchemy a persistir el valor \texttt{None} de forma explícita.

\subsubsection*{Controlador}

Expone \texttt{GET /products/}, \texttt{GET /products/\{id\}}, \texttt{POST /products/}, \texttt{PATCH /products/\{id\}} y \texttt{DELETE /products/\{id\}} sin restricciones adicionales de guards en este módulo.

% -------------------------------------------------------
\subsection{Módulo de Dispositivos}\label{sec:backend_dispositivos}

El módulo de \texttt{devices} administra las unidades físicas del inventario, su estado actual y el historial de cambios de estado. Es el módulo con mayor cantidad de métodos personalizados del proyecto.

\subsubsection*{DTOs}

\begin{itemize}
	\item \texttt{DeviceReadDTO}: excluye \texttt{loan\_request\_items} y \texttt{status\_logs}.
	\item \texttt{DeviceCreateDTO}: excluye \texttt{id}, \texttt{created\_at}, relaciones inversas y objetos de relación (\texttt{product}, \texttt{status}, \texttt{ubication}).
	\item \texttt{DeviceUpdateDTO}: igual que \texttt{CreateDTO} más \texttt{status\_id}, ya que el estado se gestiona por un endpoint dedicado.
	\item \texttt{DeviceStatusChangeDTO}: dataclass simple con un único campo \texttt{status\_id}, utilizado exclusivamente por el endpoint de cambio de estado.
	\item \texttt{DeviceStatusLogReadDTO}: expone el historial de estados excluyendo la relación inversa con \texttt{device}.
\end{itemize}

\subsubsection*{Repositorio}

\texttt{DeviceRepository} define los siguientes métodos personalizados, además de los genéricos heredados:

\begin{itemize}
	\item \texttt{list\_all\_with\_relations()}: lista todos los dispositivos con \texttt{product}, \texttt{status} y \texttt{ubication} cargados mediante \texttt{selectinload}, evitando consultas N+1 al serializar.
	\item \texttt{list\_by\_product(product\_id)}: filtra dispositivos por producto, útil para mostrar las unidades físicas de un ítem específico en el frontend.
	\item \texttt{list\_available()}: filtra dispositivos cuyo \texttt{status.name == "disponible"} mediante un \texttt{JOIN}, exponiendo solo los ítems prestables al momento de crear una solicitud.
	\item \texttt{get\_with\_relations(device\_id)}: obtiene un dispositivo por ID con todas sus relaciones cargadas.
	\item \texttt{change\_status(device\_id, status\_id, user\_id)}: actualiza el \texttt{status\_id} del dispositivo y genera automáticamente un registro en \texttt{device\_status\_logs} con el usuario que realizó el cambio y la marca temporal. Este método es el \textbf{núcleo del historial} del inventario.
\end{itemize}

\texttt{DeviceStatusLogRepository} implementa el método \\ \texttt{get\_history\_by\_device(device\_id)}, que retorna el historial de estados de un dispositivo ordenado cronológicamente.

\subsubsection*{Controlador}

\texttt{DeviceController} expone los siguientes endpoints:

\begin{itemize}
	\item \texttt{GET /devices/}: lista todos los dispositivos; acepta el parámetro de query opcional \texttt{?product\_id=} para filtrar por producto, delegando en \texttt{list\_by\_product} o \texttt{list\_all\_with\_relations} según corresponda.
	\item \texttt{GET /devices/available}: lista solo los dispositivos disponibles para préstamo.
	\item \texttt{GET /devices/\{device\_id\}}: obtiene un dispositivo por ID con sus relaciones.
	\item \texttt{POST /devices/}: registra un nuevo dispositivo en el inventario.
	\item \texttt{PATCH /devices/\{device\_id\}}: actualiza campos generales del dispositivo, excluyendo el estado.
	\item \texttt{PATCH /devices/\{device\_id\}/status}: cambia el estado del dispositivo y registra el cambio en el historial. Requiere el usuario autenticado (\texttt{request.user.id}) para atribuir el cambio.
	\item \texttt{GET /devices/\{device\_id\}/history}: retorna el historial de cambios de estado del dispositivo.
	\item \texttt{DELETE /devices/\{device\_id\}}: elimina un dispositivo; gracias al \texttt{ON DELETE CASCADE} definido en las migraciones, sus \texttt{loan\_request\_items} asociados se eliminan en cascada.
\end{itemize}

% -------------------------------------------------------
\subsection{Módulo de Préstamos}\label{sec:backend_prestamos}

El módulo de \texttt{loans} modela el ciclo de vida completo de las solicitudes de préstamo. Es el módulo de mayor complejidad funcional, pues sus operaciones involucran múltiples tablas y actualizaciones de estado coordinadas.

\subsubsection*{DTOs}

\texttt{LoanRequestReadDTO} expone la solicitud con sus relaciones anidadas: 
\texttt{user}, \texttt{status}, \texttt{loan\_request\_items} con sus \texttt{device} 
y el \texttt{product} de cada uno. Dado que estas relaciones son bidireccionales en el 
modelo ORM, sin una exclusión explícita el serializador intentaría seguir las referencias 
en ambas direcciones indefinidamente: desde \texttt{LoanRequest} hacia \texttt{Device}, 
y desde \texttt{Device} de vuelta hacia \texttt{LoanRequest} a través de 
\texttt{loan\_request\_items}, generando un ciclo infinito que termina en error. Por 
ello, el DTO excluye explícitamente campos como \texttt{loan\_request\_items.loan\_request}, 
\texttt{loan\_request\_items.device.loan\_request\_items} y otros que cierran el ciclo, 
además de ocultar datos sensibles como \texttt{user.password}.

\texttt{LoanRequestCreateDTO}, en cambio, se implementa como un \texttt{dataclass} en 
lugar de un \texttt{SQLAlchemyDTO}. Esto se debe a que el payload de creación no 
corresponde directamente a la estructura de ninguna tabla: el cliente envía una lista de 
\texttt{device\_ids}, un \texttt{reason} y una \texttt{estimated\_return\_date}, mientras 
que el repositorio se encarga de construir internamente tanto la \texttt{LoanRequest} 
como los \texttt{LoanRequestItem} asociados. Forzar este payload dentro de un 
\texttt{SQLAlchemyDTO} implicaría exponer la estructura interna del modelo ORM al cliente, 
acoplando innecesariamente la API al esquema de la base de datos~\cite{leapcell_sqlalchemy_dataclasses}. El \texttt{dataclass} 
permite definir un contrato de entrada limpio, independiente del modelo, que el repositorio 
traduce a las operaciones de persistencia correspondientes.~\cite{leapcell_decoupling_dto}

\subsubsection*{Repositorio}

\texttt{LoanRequestRepository} implementa los siguientes métodos personalizados:

\begin{itemize}
	\item \texttt{list\_with\_relations()}: lista todas las solicitudes con sus relaciones cargadas, ordenadas por \texttt{id} descendente. \textbf{Utilizado por administradores}.
	\item \texttt{list\_mine(user\_id)}: filtra las solicitudes del usuario autenticado. Permite que cada usuario vea solo sus propios préstamos.
	\item \texttt{get\_with\_relations(loan\_id)}: obtiene una solicitud por ID con todas sus relaciones.
	\item \texttt{create\_with\_items(user\_id, device\_ids, reason, estimated\_return\_date)}: crea una solicitud en estado \texttt{`pendiente'} y asocia en una sola transacción los \texttt{LoanRequestItem} correspondientes a cada \texttt{device\_id} recibido. Utiliza \texttt{session.flush()} para obtener el \texttt{id} de la solicitud antes del commit.
	\item \texttt{approve(loan\_id)}: cambia el estado de la solicitud a \texttt{`aprobado'}.
	\item \texttt{reject(loan\_id)}: cambia el estado a \texttt{`rechazado'}.
	\item \texttt{deliver(loan\_id)}: marca la solicitud como \texttt{`prestado'}, registra la \texttt{delivery\_date} y actualiza el \texttt{status\_id} de cada dispositivo asociado a \texttt{`prestado'}.
	\item \texttt{register\_return(loan\_id)}: marca la solicitud como \texttt{`devuelto'}, registra la \texttt{actual\_return\_date} y restaura el \texttt{status\_id} de cada dispositivo a \texttt{`disponible'}.
\end{itemize}

\subsubsection*{Controlador}

\texttt{LoanRequestController} expone los siguientes endpoints:

\begin{itemize}
	\item \texttt{GET /loans/}: lista todas las solicitudes; protegido por \texttt{admin\_guard}, solo accesible para administradores.
	\item \texttt{GET /loans/me}: lista las solicitudes del usuario autenticado.
	\item \texttt{GET /loans/\{loan\_id\}}: detalle de una solicitud con sus ítems y dispositivos.
	\item \texttt{POST /loans/}: crea una solicitud de préstamo; el \texttt{user\_id} se extrae del token (\texttt{request.user.id}), no del payload, para garantizar que el usuario no pueda crear solicitudes en nombre de otros.
	\item \texttt{PATCH /loans/\{loan\_id\}/approve}: aprueba una solicitud; protegido por \texttt{admin\_guard}.
	\item \texttt{PATCH /loans/\{loan\_id\}/reject}: rechaza una solicitud; protegido por \texttt{admin\_guard}.
	\item \texttt{PATCH /loans/\{loan\_id\}/deliver}: registra la entrega; protegido por \texttt{admin\_guard}.
	\item \texttt{PATCH /loans/\{loan\_id\}/return}: registra la devolución; protegido por \texttt{admin\_guard}.
	\item \texttt{DELETE /loans/\{loan\_id\}}: elimina una solicitud; gracias al \texttt{ON DELETE CASCADE}, sus \texttt{loan\_request\_items} se eliminan en cascada.
\end{itemize}